\documentclass[11pt,a4paper]{article}
\usepackage[T1]{fontenc}
\usepackage[utf8]{inputenc}
\usepackage[main=italian,provide=*]{babel}
\usepackage{amsmath,amssymb}
\usepackage{mathtools}
\usepackage{xcolor}
\usepackage[a4paper,margin=2.7cm]{geometry}
\usepackage{caption}
\usepackage{booktabs}
\usepackage{enumitem}
\usepackage{placeins}
\usepackage[colorlinks=true,linkcolor=black,citecolor=black,urlcolor=blue!55!black]{hyperref}
\newcommand{\ket}[1]{\lvert #1 \rangle}
\newcommand{\bra}[1]{\langle #1 \rvert}

\title{Guida d'uso --- pacchetto riproducibile\\Parte 2 (rumore), trimero ad anello}
\author{Tesi triennale, Samuele --- Relatore: Prof.\ Alessandro Chiesa}
\date{}

\begin{document}
\sloppy
\maketitle

\tableofcontents

\section{Struttura del pacchetto}

\begin{verbatim}
codice/            moduli .py (libreria, mai eseguiti direttamente)
dati/              output .npz degli script di generazione dati
figure/            figure finali, PDF+PNG
genera_figure/     uno script per figura
01_..., 07_...     script di dati/validazione, numerati nell'ordine d'uso
esegui_tutto.py    orchestratore: dati -> validazioni -> figure
\end{verbatim}

Comando singolo per rigenerare tutto da zero:
\begin{verbatim}
python3 esegui_tutto.py
\end{verbatim}
Richiede: \texttt{qiskit}, \texttt{qiskit-aer}, \texttt{numpy}, \texttt{scipy}, \texttt{matplotlib}.

\section{Dipendenza incrociata con Parte 1, dichiarata}

Cinque moduli in \texttt{codice/} sono copie bit-per-bit di file del
Project che appartengono concettualmente a \textbf{Parte 1} (sistema
chiuso), non a Parte 2: \texttt{vqe\_w2q6\_\allowbreak trimero\_anello.py},
\texttt{trimer\_ring\_exact.py}, \texttt{trotter\_\allowbreak trimero\_anello.py},
\texttt{circuito\_\allowbreak correlazioni\_\allowbreak trimero\_anello.py},
\texttt{noise\_model\_dimero.py} (quest'ultimo dal dimero, non
dall'anello). Riusati qui invece di duplicarne la logica --- vedi i
docstring dei moduli di Parte 2 che li importano per il dettaglio di
cosa viene preso da ciascuno.

\section{Le due onde di dipendenza reale}

\textbf{Onda 1} (nessun ingresso): solo \texttt{01}. \textbf{Onda 2}
(dipendono tutti da \texttt{01}, indipendenti fra loro):
\texttt{02}--\texttt{07}. \textbf{Onda 3} (dipendono dai \texttt{.npz}
prodotti nell'Onda 2): gli script in \texttt{genera\_figure/}.

\section{Script 01--07}

\subsection{01\_genera\_parametri\_vqe.py}
\textbf{Ingresso}: nessuno (calcola tutto da zero: Hamiltoniana +
ottimizzazione multistart). \textbf{Uscita}:
\texttt{dati/w2q6\_params\_optimal.npz} (campi: \texttt{params},
\texttt{E\_vqe}, \texttt{E\_exact}, \texttt{fidelity}). \textbf{Chi lo
consuma}: tutti gli script \texttt{02}--\texttt{07}.
\begin{verbatim}
python3 01_genera_parametri_vqe.py
\end{verbatim}

\subsection{02\_valida\_modello\_rumore.py}
\textbf{Ingresso}: \texttt{dati/w2q6\_params\_optimal.npz}.
\textbf{Uscita}: nessun file, solo resoconto a schermo (tre controlli
sul modello di rumore, Stadio 1). \textbf{Chi lo consuma}: nessuno.
\begin{verbatim}
python3 02_valida_modello_rumore.py
\end{verbatim}

\subsection{03\_valida\_vqe\_dm\_rumoroso.py}
\textbf{Ingresso}: \texttt{dati/w2q6\_params\_optimal.npz}.
\textbf{Uscita}: \texttt{dati/scan\_eps2q\_\allowbreak vqe\_dm\_trimero\_anello.npz}
(Stadio 2). \textbf{Chi lo consuma}:
\texttt{genera\_figure/\allowbreak genera\_scan\_eps2q\_\allowbreak trimero\_anello.py}.
\begin{verbatim}
python3 03_valida_vqe_dm_rumoroso.py
\end{verbatim}

\subsection{04\_valida\_trotter\_rumoroso.py}
\textbf{Ingresso}: \texttt{dati/w2q6\_params\_optimal.npz}.
\textbf{Uscita}:
\texttt{dati/scan\_N\_trotter\_\allowbreak rumoroso\_trimero\_anello.npz} (Stadio 3,
entrambi gli scenari A/B, griglia $N$ fine --- vedi nota sotto).
\textbf{Chi lo consuma}:
\texttt{genera\_figure/\allowbreak genera\_trotter\_\allowbreak rumoroso\_trimero\_anello.py}.
\begin{verbatim}
python3 04_valida_trotter_rumoroso.py
\end{verbatim}

\subsection{05\_valida\_correlatori\_rumorosi.py}
\textbf{Ingresso}: \texttt{dati/w2q6\_params\_optimal.npz}.
\textbf{Uscita}: \texttt{dati/scan\_N\_\allowbreak correlatore\_trimero\_anello.npz}
(Stadio 4). \textbf{Chi lo consuma}:
\texttt{genera\_figure/\allowbreak genera\_correlatore\_\allowbreak rumoroso\_trimero\_anello.py}.
\begin{verbatim}
python3 05_valida_correlatori_rumorosi.py
\end{verbatim}

\subsection{06\_genera\_scan\_parametri.py}
\textbf{Ingresso}: \texttt{dati/w2q6\_params\_optimal.npz}.
\textbf{Uscita}:
\texttt{dati/scan\_parametri\_\allowbreak rumore\_trimero\_anello.npz} (Stadio 5,
entrambi gli scenari). \textbf{Chi lo consuma}:
\texttt{genera\_figure/\allowbreak genera\_Nstar\_vs\_eps\_\allowbreak trimero\_anello.py}.
\begin{verbatim}
python3 06_genera_scan_parametri.py
\end{verbatim}

\subsection{07\_valida\_scan\_parametri.py}
\textbf{Ingresso}: \texttt{dati/w2q6\_params\_optimal.npz}.
\textbf{Uscita}: nessun file, solo resoconto a schermo (regressione,
monotonia, invarianza da readout). \textbf{Chi lo consuma}: nessuno.
\begin{verbatim}
python3 07_valida_scan_parametri.py
\end{verbatim}

\section{Script genera\_figure/}

\subsection{genera\_depolarizzazione\_bloch.py}
\textbf{Ingresso}: nessuno (formula analitica $1-\lambda$).
\textbf{Uscita}: \texttt{figure/fig\_\allowbreak depolarizzazione\_bloch.pdf/.png}.
\begin{verbatim}
python3 genera_figure/genera_depolarizzazione_bloch.py
\end{verbatim}

\subsection{genera\_scan\_eps2q\_trimero\_anello.py}
\textbf{Ingresso}: \texttt{dati/scan\_eps2q\_\allowbreak vqe\_dm\_trimero\_anello.npz}.
\textbf{Uscita}: \texttt{figure/fig\_scan\_eps2q\_\allowbreak trimero\_anello.pdf/.png}.
\begin{verbatim}
python3 genera_figure/genera_scan_eps2q_trimero_anello.py
\end{verbatim}

\subsection{genera\_trotter\_rumoroso\_trimero\_anello.py}
\textbf{Ingresso}:
\texttt{dati/scan\_N\_trotter\_\allowbreak rumoroso\_trimero\_anello.npz}.
\textbf{Uscita}:
\texttt{figure/fig\_trotter\_\allowbreak rumoroso\_trimero\_anello.pdf/.png}.
\begin{verbatim}
python3 genera_figure/genera_trotter_rumoroso_trimero_anello.py
\end{verbatim}

\subsection{genera\_correlatore\_rumoroso\_trimero\_anello.py}
\textbf{Ingresso}: \texttt{dati/scan\_N\_\allowbreak correlatore\_trimero\_anello.npz}.
\textbf{Uscita}:
\texttt{figure/fig\_correlatore\_\allowbreak rumoroso\_trimero\_anello.pdf/.png}.
\begin{verbatim}
python3 genera_figure/genera_correlatore_rumoroso_trimero_anello.py
\end{verbatim}

\subsection{genera\_Nstar\_vs\_eps\_trimero\_anello.py}
\textbf{Ingresso}: \texttt{dati/scan\_parametri\_\allowbreak rumore\_trimero\_anello.npz}.
\textbf{Uscita}: \texttt{figure/fig\_Nstar\_vs\_eps\_\allowbreak trimero\_anello.pdf/.png}.
\begin{verbatim}
python3 genera_figure/genera_Nstar_vs_eps_trimero_anello.py
\end{verbatim}

\section{Nota: correzione della griglia \texorpdfstring{$N$}{N} (Stadio 3)}

Una prima versione dello script 04 usava una griglia $N$ troppo grezza
($\{1,2,5,10,15,\ldots\}$), che saltava $N=3$ e $N=4$ e riportava
$N^*_A=2$ invece del vero $3$. Corretto con una griglia continua
$N=1,\ldots,20$ (poi diradata) --- lo script in questo pacchetto usa già
la versione corretta. Vedi
\texttt{risultati\_scan\_parametri\_rumore\_trimero\_anello.tex},
Sez.\ ``Correzione'', per i dettagli.

\section{Sequenza completa (copia-incolla)}

\begin{verbatim}
python3 01_genera_parametri_vqe.py
python3 02_valida_modello_rumore.py
python3 03_valida_vqe_dm_rumoroso.py
python3 04_valida_trotter_rumoroso.py
python3 05_valida_correlatori_rumorosi.py
python3 06_genera_scan_parametri.py
python3 07_valida_scan_parametri.py
python3 genera_figure/genera_depolarizzazione_bloch.py
python3 genera_figure/genera_scan_eps2q_trimero_anello.py
python3 genera_figure/genera_trotter_rumoroso_trimero_anello.py
python3 genera_figure/genera_correlatore_rumoroso_trimero_anello.py
python3 genera_figure/genera_Nstar_vs_eps_trimero_anello.py
\end{verbatim}
Equivalente, in un solo comando: \texttt{python3 esegui\_tutto.py}.

\section{Verifica end-to-end effettuata}

\texttt{dati/} e \texttt{figure/} svuotate, \texttt{esegui\_tutto.py}
rieseguito da zero: \texttt{exit code 0}, nessun \texttt{FALLITO} n\'e
\texttt{Traceback} nel log. Numeri rigenerati confrontati con quelli nei
documenti di risultati: $E_\text{VQE}=-5.534370017393404$,
$N^*_A=3$ ($F=0.700335$), $N^*_B=9$ ($F=0.381896$), $N^*$
correlatore$=3$ ($|C|=0.505$) --- tutti coincidenti. Tutte le cinque
figure verificate vettoriali (\texttt{pdffonts}: zero Type 3;
\texttt{pdfimages -list}: zero immagini raster incorporate).

\end{document}
